%% ma: L11-B16-0088
%% lop: 11
%% chuong: 5
%% bai: 16
%% dang: 11.16.1
%% loai: TL
%% muc: VD
%% dap_an: "-"
%% nguon: "Ảnh giáo viên gửi 25/11/2026: Giới hạn hàm số dạng 0/0 (3 bài, có ĐS), Bài 2 câu 6"
%% kien_thuc: "Lớp 11 Bài 16 (giới hạn của hàm số tại một điểm); chia đa thức theo sơ đồ Horner, phân tích đa thức thành nhân tử"
%% trang_thai: chinh-thuc
\begin{ex}
Tính giới hạn $A=\lim\limits_{x\to3}\dfrac{x^3-5x^2+3x+9}{x^4-8x^2-9}$.
\loigiai{$x^3-5x^2+3x+9=(x-3)^2(x+1)$, $x^4-8x^2-9=(x-3)(x+3)(x^2+1)$. Biểu thức bằng $\dfrac{(x-3)(x+1)}{(x+3)(x^2+1)}$. Vậy $A=0$.}
\end{ex}
